\beginchapter Chapter 17. More about Math

Another thing mathematicians like to do is make fractions---and they
like to build symbols up on top of each other in a variety of different ways:
\begindisplay
$\displaystyle
{1\over2}\qquad{\rm and}\qquad{n+1\over3}\qquad{\rm and}\qquad
{n+1\choose3}\qquad{\rm and}\qquad\sum_{n=1}^3 Z_n^2\,.$
\enddisplay
You can get these four formulas as displayed equations by typing
`|$$1\over2$$|' and
`|$$n+1\over3$$|' and
`|$$n+1\choose3$$|' and
`|$$\sum_{n=1}^3 Z_n^2$$|';
we shall study the simple rules for such constructions in this chapter.
^^|\sum|^^|\choose|

First let's look at ^{fractions}, which use the `^|\over|' notation. The
control sequence |\over| applies to everything in the formula unless you
use braces to enclose it in a specific subformula; in the latter
^^{stacked fractions, see over}
case, |\over| applies to everything in that subformula.
\begindisplaymathdemo
\it Input&\it Output\cr
\noalign{\vskip-3pt}
|$$x+y^2\over k+1$$|&x+y^2\over k+1\cr
\noalign{\vskip2pt}
|$${x+y^2\over k}+1$$|&{x+y^2\over k}+1\cr
\noalign{\vskip-1pt}
|$$x+{y^2\over k}+1$$|&x+{y^2\over k}+1\cr
\noalign{\vskip-1pt}
|$$x+{y^2\over k+1}$$|&x+{y^2\over k+1}\cr
\noalign{\vskip-3pt}
|$$x+y^{2\over k+1}$$|&x+y^{2\over k+1}\cr
\endmathdemo
You aren't allowed to use |\over| twice in the same subformula; instead of
typing something like
`|a \over b \over 2|', you must specify what goes over what:
\begindisplaymathdemo
\noalign{\vskip3pt}
|$${a\over b}\over 2$$|&{a\over b}\over 2\cr
|$$a\over{b\over 2}$$|&a\over{b\over 2}\cr
\endmathdemo
Unfortunately, both of these alternatives look pretty awful. Mathematicians
tend to ``overuse'' |\over| when they first begin to typeset their own work
on a system like \TeX. A good typist or copy editor will convert fractions
to a ``^{slashed form},'' whenever a built-up construction would be too
small or too crowded. For example, the last two cases should be treated
as follows:
\begindisplaymathdemo
\noalign{\vskip3pt}
|$$a/b \over 2$$|&a/b \over 2\cr
|$$a \over b/2$$|&a \over b/2\cr
\endmathdemo
Conversion to slashed form takes a little bit of mathematical knowhow, since
^{parentheses} sometimes need to be inserted in order to preserve the meaning
of the formula. Besides substituting `|/|' for~`|\over|', the two parts
of the fraction should be put in parentheses unless they are single
symbols; for example, $a\over b$~becomes simply~$a/b$, but
$a+1\over b$ becomes $(a+1)/b$, and $a+1\over b+1$ becomes
${(a+1)/(b+1)}$. Furthermore, the entire fraction should generally
be enclosed in parentheses if it appears next to something else;
for example, ${a\over b}x$ becomes $(a/b)x$. If you are a typist without
mathematical training, it's best to ask the author of the manuscript
for help, in doubtful cases; you might also tactfully suggest that
unsightly fractions be avoided altogether in future manuscripts.

\exercise What's a better way to render the formula $x+y^{2\over k+1}$?
\answer $x+y^{2/(k+1)}$\quad(|$x+y^{2/(k+1)}$|).

\exercise Convert `${a+1\over b+1}x$' to slashed form.
\answer $((a+1)/(b+1))x$\quad(|$((a+1)/(b+1))x$|).

\exercise What surprise did B. L. ^{User} get when he typed `|$$x = (y^2\over
k+1)$$|'\thinspace?
\answer He got the displayed formula$$x=(y^2\over k+1)$$ because he forgot
that an unconfined |\over| applies to everything.  \ (He should probably
have typed `|$$x=\left(y^2\over k+1\right)$$|', using ideas that will be
presented later in this chapter; this not only makes the parentheses
larger, it keeps the `$x=$' out of the fraction, because |\left| and
|\right| introduce subformulas.)

\def\cents{\hbox{\rm\rlap/c}}
\exercise How can you make `$7{1\over2}\cents$'? \ (Assume that
the control sequence |\cents| yields~`$\cents$'.)^^{money}^^{cents}
\answer `|$7{1\over2}\cents$|' or `|7$1\over2$\cents|'. \ (Incidentally,
the definition used here was |\def\cents{\hbox{\rm\rlap/c}}|.)
^^|\rlap|^^|\cents|

The examples above show that letters and other symbols sometimes get
smaller when they appear in fractions, just as they get smaller when they
are used as exponents. It's about time that we studied \TeX's method for
choosing the sizes of things. \TeX\ actually has eight different
^{styles} in which it can treat formulas, namely
$$\halign{\indent#\hfil\quad&#\hfil\cr
display style&(for formulas displayed on lines by themselves)\cr
text style&(for formulas embedded in the text)\cr
script style&(for formulas used as superscripts or subscripts)\cr
scriptscript style&(for second-order superscripts or subscripts)\cr}$$
^^{display style}^^{text style}^^{script style}^^{scriptscript style}
and four other ``^{cramped}'' styles that are almost the same except that
exponents aren't raised quite so much. For brevity we shall refer to the
eight styles as
\begindisplay
$\displaystyle D,\ D',\ T,\ T',\ S,\ S',\ \SS,\ \SS',$
\enddisplay
where $D$ is display style, $D'$ is cramped display style, $T$~is text style,
etc. \TeX\ also uses three different ^{sizes of type for mathematics};
they are called ^{text size}, ^{script size}, and ^{scriptscript size}.

The normal way to typeset a formula with \TeX\ is to enclose it in dollar
signs |$|$\,\ldots\,$|$|; this yields the formula in text style
(style~$T$). Or you can enclose it in double dollar signs |$$|$\,\ldots\,$|$$|;
this displays the formula in display style (style~$D$). The subformulas of
a formula might, of course, be in different styles. Once you know
the style, you can determine the size of type that \TeX\ will use:
$$\everycr{\noalign{\penalty10000}}
\halign{\indent#\hfil\qquad&#\hfil&\quad#\llap(like this)\hfil\cr
If a letter is in style&then it will be set in\cr
\noalign{\vskip 2pt}
$D,D',T,T'$&text size&\cr
$S,S'$&script size&\sevenrm\cr
$\SS,\SS'$&scriptscript size&\fiverm\cr}$$
There is no ``$\it SSS$'' style or ``scriptscriptscript'' size; such tiny
symbols would be even less readable than the scriptscript ones. Therefore
\TeX\ stays with scriptscript size as the minimum:
$$\halign{\indent\hbox to 1.3in{#\hfil}&\hbox to 1.2in{#\hfil}&#\hfil\cr
In a formula&the superscript&and the subscript\cr
of style&style is&style is\cr
\noalign{\vskip 2pt}
$D,T$&$S$&$S'$\cr
$D',T'$&$S'$&$S'$\cr
$S,\SS$&$\SS$&$\SS'$\cr
$S',\SS'$&$\SS'$&$\SS'$\cr}$$
For example, if |x^{a_b}| is to be typeset in style $D$, then |a_b| will
be set in style~$S$, and {\tt b}~in style~$\SS'$; the result is
`$\displaystyle x^{a_b}$'.

So far we haven't seen any difference between styles $D$ and $T$. Actually
there is a slight difference in the positioning of exponents, although
script size is used in each case: You get
$\displaystyle x^2$~in $D$~style and $x^2$~in $T$~style and \vbox to 0pt{
\vss\hbox{$\displaystyle{\atop x^2}$}\kern0pt}~in $D'$ or $T'$~style---do
you see the difference? But there is a big distinction between $D$ style and
$T$ style when it comes to fractions:
$$\halign{\indent\hbox to 1.3in{#\hfil}&\hbox to 1.2in{#\hfil}&#\hfil\cr
In a formula&the style of the&and the style of the\cr
$\alpha$|\over|$\,\beta$ of style&numerator $\alpha$ is&denominator
$\beta$ is\cr
\noalign{\vskip 2pt}
$D$&$T$&$T'$\cr
$D'$&$T'$&$T'$\cr
$T$&$S$&$S'$\cr
$T'$&$S'$&$S'$\cr
$S,\SS$&$\SS$&$\SS'$\cr
$S',\SS'$&$\SS'$&$\SS'$\cr}$$
^^{numerator}^^{denominator}
Thus if you type `|$1\over2$|' (in a text) you get $1\over2$, namely style
$S$ over style~$S'$; but if you type
`|$$1\over2$$|' you get $$1\over2$$ (a displayed formula), which is style
$T$ over style $T'$.

\danger While we're at it, we might as well finish the style rules:
^|\underline| does not change the style. ^{Math accents}, and the operations
^|\sqrt| and ^|\overline|, change uncramped styles to their cramped
counterparts; for example, $D$ changes to $D'$, but $D'$ stays as it was.

\dangerexercise State the style and size of each part of the formula
$\displaystyle \sqrt{p_2^{e'}}$, assuming that the formula itself is in
style~$D$.
\answer Style $D'$ is used for the subformula $p_2^{e'}$, hence style~$S'$
is used for the superscript~$e'$ and the subscript~2, and style~$\SS'$
is used for the supersuperscript prime. The square root sign and the $p$
appear in text size; the 2 and the~$e$ appear in script size; and the
$\prime$ is in scriptscript size.

Suppose you don't like the style that \TeX\ selects by its automatic style
rules.  Then you can specify the style you want by typing ^|\displaystyle|
or ^|\textstyle| or ^|\scriptstyle| or ^|\scriptscriptstyle|; the style
that you select will apply until the end of the formula or subformula, or
until you select another style.  For example,
`|$$n+\scriptstyle n+\scriptscriptstyle n.$$|' produces the display
$$n+\scriptstyle n+\scriptscriptstyle n.$$
This is a rather silly example, but it does show
that the plus signs get smaller too, as the style changes. \TeX\ puts no
space around + signs in script styles.

Here's a more useful example of style changes: Sometimes you need to
typeset a ``^{continued fraction}'' made up of many other fractions,
all of which are supposed to be in display style:
$$a_0+{1\over\displaystyle a_1+
          {\strut 1\over\displaystyle a_2+
            {\strut 1\over\displaystyle a_3+
              {\strut 1\over a_4}}}}$$
In order to get this effect, the idea is to type
\begintt
$$a_0+{1\over\displaystyle a_1+
        {\strut 1\over\displaystyle a_2+
          {\strut 1\over\displaystyle a_3+
            {\strut 1\over a_4}}}}$$
\endtt
(The control sequence ^|\strut| has been used to make the denominators
taller; this is a refinement that will be discussed in
Chapter~18. Our concern now is with the style commands.) \
Without the appearances of\/ |\strut| and |\displaystyle| in this formula,
the result would be completely different:
$$a_0+{1\over a_1+{1\over
      a_2+{1\over a_3+{1\over a_4}}}}$$

\danger These examples show that the numerator and denominator of a fraction
are generally centered with respect to each other. If you prefer to have
the numerator or denominator appear ^{flush left}, put `^|\hfill|' after
it; or if you prefer ^{flush right}, put `|\hfill|' at the left. For
example, if the first three appearances of `|1\over|' in the previous
example are replaced by `|1\hfill\over|', you get the display
$$a_0+{1\hfill\over\displaystyle a_1+
          {\strut1\hfill\over\displaystyle a_2+
            {\strut1\hfill\over\displaystyle a_3+
              {\strut1\over a_4}}}}$$
(a format for continued fractions that many authors prefer). This works
because |\hfill| stretches at a faster rate than the glue that is
actually used internally by \TeX\ when it centers the numerators
and denominators.

\TeX\ has another operation `^|\atop|', which is like |\over| except that
it leaves out the fraction line:
\begindisplaymathdemo
|$$x\atop y+2$$|&x\atop y+2\cr
\endmathdemo
The plain \TeX\ format in Appendix B also defines `^|\choose|', which is
like |\atop| but it encloses the result in parentheses:
\begindisplaymathdemo
|$$n\choose k$$|&n\choose k\cr
\endmathdemo
It is called |\choose| because it's
a common notation for the so-called ^{binomial coefficient}
that tells how many ways there are to choose $k$~things out of $n$~things.

You can't mix |\over| and |\atop| and |\choose| with each other.
For example, `|$$n \choose k \over 2$$|' is illegal; you must use
grouping, to get either `|$${n\choose k}\over2$$|' or
`|$$n\choose{k\over2}$$|', i.e.,
\begindisplay
$\displaystyle{{n\choose k}\over2}\qquad{\rm or}\qquad {n\choose{k\over2}}.$
\enddisplay
The latter formula, incidentally, would look better as
`|$$n\choose k/2$$|' or `|$$n\choose{1\over2}k$$|', yielding
\begindisplay
$\displaystyle{n\choose k/2}\qquad{\rm or}\qquad{n\choose{1\over2}k}.$
\enddisplay

\medskip
\exercise As alternatives to $\displaystyle{{n\choose k}\over2}$,
discuss how you could obtain the two displays
\begindisplay\abovedisplayskip=0pt\belowdisplayskip=0pt
$\displaystyle
{1\over2}{n\choose k}
\qquad{\rm and}\qquad
{\displaystyle{n\choose k}\over2}.$
\enddisplay
\answer |$${1\over2}{n\choose k}$$|;
|$$\displaystyle{n\choose k}\over2$$|.
All of these braces are necessary.

\bigbreak
\exercise Explain how to specify the displayed formula
$${p \choose 2}x^2 y^{p-2} - {1 \over 1-x}{1 \over 1-x^2}.$$
\answer |$${p \choose 2} x^2 y^{p-2} - {1 \over 1-x}{1 \over 1-x^2}.$$|

\danger \TeX\ has a generalized version of\/ |\over| and |\atop| in which you
specify the exact thickness of the line rule by typing
`^|\above|\<dimen>'. For example,
\begintt
$$\displaystyle{a\over b}\above1pt\displaystyle{c\over d}$$
\endtt
will produce a ^{compound fraction} with a heavier ($1\pt$ thick) rule as
its main bar:
$${\displaystyle{a\over b}\above 1pt\displaystyle{c\over d}}.$$
This sort of thing occurs primarily in textbooks on elementary mathematics.

\goodbreak
Mathematicians often use the sign $\sum$ to stand for ``^{summation}''
and the sign $\int$ to stand for ``^{integration}.'' If you're a typist but not
a mathematician, all you need to remember is that ^|\sum| stands for
$\sum$ and ^|\int| for $\int$; these abbreviations appear in Appendix~F
together with all the other symbols, in case you forget. Symbols like
$\sum$ and $\int$ (and a few others like $\bigcup$ and $\prod$ and $\oint$
and~$\bigotimes$, all listed in Appendix~F) are called {\sl ^{large operators}},
^^{collective signs, see large operators} ^^{sigma signs, see sum}
and you type them just as you type ordinary symbols or letters. The
difference is that \TeX\ will choose a {\sl larger\/} large operator in
display style than it will in text style. For example,
$$\halign{\indent#\hfil\qquad yields\qquad&$#\hfil$\qquad&#\hfil\cr
|$\sum x_n$|&\sum x_n&($T$ style)\cr
\noalign{\vskip3pt}
|$$\sum x_n$$|&\displaystyle\sum x_n&($D$ style).\cr}$$

A displayed |\sum| usually occurs with ``^{limits},'' i.e., with
subformulas that are to appear above and below it. You type limits just
as if they were superscripts and subscripts; for example, if you want
$$\sum_{n=1}^m$$
you type either `|$$\sum_{n=1}^m$$|' or `|$$\sum^m_{n=1}$$|'.  According
to the normal conventions of mathematical typesetting, \TeX\ will change
this to `$\sum_{n=1}^m$' (i.e., without limits) if it occurs in text
style rather than in display style.

Integrations are slightly different from summations, in that the superscripts
and subscripts are not set as limits even in display style:
$$\halign{\indent\hbox to2.3in{#\hfil}\hbox to.6in{yields\hfil}&
  $#\hfil$\qquad&#\hfil\cr
|$\int_{-\infty}^{+\infty}$|&\int_{-\infty}^{+\infty}&($T$ style)\cr
\noalign{\vskip3pt}
|$$\int_{-\infty}^{+\infty}$$|&\displaystyle\int_{-\infty}^{+\infty}&
  ($D$ style).\cr}$$

\danger Some printers prefer to set limits above and below $\int$ signs;
this takes more space on the page, but it
gives a better appearance if the subformulas are complex, because it
keeps them out of the way of the rest of the formula. Similarly, limits
are occasionally desirable in text style or script style; but some
printers prefer not to set limits on displayed $\sum$ signs. You can change
\TeX's convention by simply typing `^|\limits|' or `^|\nolimits|' immediately
after the large operator.
For example,
$$\halign{\indent\hbox to2.3in{#\hfil}\hbox to.6in{yields\hfil}&
  $\displaystyle{#}$\hfil\cr
|$$\int\limits_0^{\pi\over2}$$|&\int\limits_0^{\pi\over2}\cr
\noalign{\vskip 4pt}
|$$\sum\nolimits_{n=1}^m$$|&\sum\nolimits_{n=1}^m\cr}$$

\ddanger If you say `|\nolimits\limits|' (presumably because some macro
like |\int| specifies |\nolimits|, but you do want them), the last word
takes precedence.  There's also a command `^|\displaylimits|' that can be
used to restore \TeX's normal conventions; i.e., the limits will be
displayed only in styles $D$ and $D'$.

\danger Sometimes you need to put two or more rows of limits under a large
operator; you can do this with `^|\atop|'. For example, if you want
the displayed formula
$$\sum_{\scriptstyle0\le i\le m\atop\scriptstyle0<j<n}P(i,j)$$
the correct way to type it is
\begintt
$$\sum_{\scriptstyle0\le i\le m\atop\scriptstyle0<j<n}P(i,j)$$
\endtt
(perhaps with a few more spaces to make it look nicer in the manuscript
file). The instruction `^|\scriptstyle|' was necessary here,
twice---otherwise the lines `$0\le i\le m$' and `$0<j<n$' would have been in
scriptscript size, which is too small. This is another instance of a rare
case where \TeX's automatic style rules need to be overruled.

\exercise How would you type the displayed formula $\displaystyle
\sum_{i=1}^p\sum_{j=1}^q\sum_{k=1}^ra_{ij}b_{jk}c_{ki}$\enspace?
\answer |$$\sum_{i=1}^p\sum_{j=1}^q\sum_{k=1}^ra_{ij}b_{jk}c_{ki}$$|.

\dangerexercise And how would you handle $\displaystyle
\sum_{{\scriptstyle1\le i\le p\atop\scriptstyle1\le j\le q}
    \atop\scriptstyle1\le k\le r}a_{ij}b_{jk}c_{ki}$\enspace?
\answer |$$\sum_{{\scriptstyle 1\le i\le p \atop \scriptstyle 1\le j\le q}
    \atop \scriptstyle 1\le k\le r} a_{ij} b_{jk} c_{ki}$$|.

Since mathematical formulas can get horribly large, \TeX\ has to have some
way to make ever-larger symbols. For example, if you type
\begintt
$$\sqrt{1+\sqrt{1+\sqrt{1+
            \sqrt{1+\sqrt{1+\sqrt{1+\sqrt{1+x}}}}}}}$$
\endtt
the result shows a variety of available ^{square-root signs}:
\begindisplay
$\displaystyle\sqrt{1+\sqrt{1+\sqrt{1+
            \sqrt{1+\sqrt{1+\sqrt{1+\sqrt{1+x}}}}}}}$
\enddisplay
The three largest signs here are all essentially the same, except for a
vertical segment `\vbox{\hbox{\tenex\char'165}\vss}' that gets repeated as
often as necessary to reach the desired size; but the smaller signs are
distinct characters found in \TeX's math fonts.

A similar thing happens with parentheses and other so-called
``^{delimiter}'' symbols.  For example, here are some of the different sizes of
^^{fences, see delimiters}
^{parentheses} and ^{braces} that plain \TeX\ might use in formulas:
\begindisplay
$\displaystyle
\left(\vbox to 27pt{}\left(\vbox to 24pt{}\left(\vbox to 21pt{}
\Biggl(\biggl(\Bigl(\bigl(({\scriptstyle({\scriptscriptstyle(\hskip3pt
)})})\bigr)\Bigr)\biggr)\Biggr)\right)\right)\right)
\left\{\vbox to 27pt{}\left\{\vbox to 24pt{}\left\{\vbox to 21pt{}
\Biggl\{\biggl\{\Bigl\{\bigl\{\{{\scriptstyle\{{\scriptscriptstyle\{\hskip3pt
\}}\}}\}\bigr\}\Bigr\}\biggr\}\Biggr\}\right\}\right\}\right\}$
\enddisplay
The three largest pairs in each case are made with repeatable extensions,
so they can become as large as necessary.
^^{pieces of symbols}

Delimiters are important to mathematicians, because they provide good
visual clues to the underlying structure of complex expressions; they delimit
the boundaries of individual subformulas. Here is a list of the 22~basic
delimiters provided by plain \TeX:
\begindisplay
\it Input&\it Delimiter\cr
\noalign{\vskip2pt}
|(|&left parenthesis: $($\cr
|)|&right parenthesis: $)$\cr
|[| or ^|\lbrack|&left bracket: $[$\cr
|]| or ^|\rbrack|&right bracket: $]$\cr
|\{| or ^|\lbrace|&left curly brace: $\{$\cr
|\}| or ^|\rbrace|&right curly brace: $\}$\cr
^|\lfloor|&left floor bracket: $\lfloor$\cr
^|\rfloor|&right floor bracket: $\rfloor$\cr
^|\lceil|&left ceiling bracket: $\lceil$\cr
^|\rceil|&right ceiling bracket: $\rceil$\cr
^|\langle|&left angle bracket: $\langle$\cr
^|\rangle|&right angle bracket: $\rangle$\cr
|/|&slash: $/$\cr
^|\backslash|&reverse slash: $\backslash$\cr
\| or ^|\vert|&vertical bar: $\vert$\cr
|\|\| or ^|\Vert|&double vertical bar: $\Vert$\cr
^|\uparrow|&upward arrow: $\uparrow$\cr
^|\Uparrow|&double upward arrow: $\Uparrow$\cr
^|\downarrow|&downward arrow: $\downarrow$\cr
^|\Downarrow|&double downward arrow: $\Downarrow$\cr
^|\updownarrow|&up-and-down arrow: $\updownarrow$\cr
^|\Updownarrow|&double up-and-down arrow: $\Updownarrow$\cr
\enddisplay
^^{bent bars, see langle, rangle} ^^{curly braces, see lbrace, rbrace}
^^{leftbracket}^^{rightbracket}^^{leftbrace}^^{rightbrace}^^{/}
In some cases, there are two ways to get the same delimiter; for example,
you can specify a left bracket by typing either `|[|' or `|\lbrack|'. The
latter alternative has been provided because the symbol `|[|' is not
readily available on all computer keyboards. Remember, however,
that you should never try to specify a left brace or right brace simply by
typing `|{|' or `|}|'; the |{| and |}| symbols are reserved for grouping.
The right way is to type `|\{|' or `|\}|' or `|\lbrace|' or `|\rbrace|'.

In order to get a slightly larger version of any of these symbols, just
precede them by `^|\bigl|' (for opening delimiters) or `^|\bigr|' (for
closing ones). This makes it easier to read formulas that contain
delimiters inside delimiters:
\beginlongmathdemo
\it Input&\it Output\cr
\noalign{\vskip2pt}
|$\bigl(x-s(x)\bigr)\bigl(y-s(y)\bigr)$|&
  \bigl(x-s(x)\bigr)\bigl(y-s(y)\bigr)\cr
|$\bigl[x-s[x]\bigr]\bigl[y-s[y]\bigr]$|&
  \bigl[x-s[x]\bigr]\bigl[y-s[y]\bigr]\cr
|$\bigl|\|| |\||x|\||-|\||y|\|| \bigr|\||$|&
  \bigl\vert\vert x\vert-\vert y\vert\bigr\vert\cr
|$\bigl\lfloor\sqrt A\bigr\rfloor$|&
  \bigl\lfloor\sqrt A\bigr\rfloor\cr
\endmathdemo
The |\big| delimiters are just enough bigger than ordinary ones so that
the difference can be perceived, yet small enough to be used in the text
of a paragraph. Here are all~22 of them, in the ordinary size and in
the |\big| size:
\begindisplay
$(\,)\,[\,]\,\{\,\}\,\lfloor\,\rfloor\,\lceil\,\rceil\,\langle\,\rangle
  \,/\,\backslash\,\vert\,\Vert\,\uparrow\,\Uparrow\,\downarrow\,\Downarrow
  \,\updownarrow\,\Updownarrow$\cr
\noalign{\smallskip}
$\bigl(\,\bigr)\,\bigl[\,\bigr]\,\bigl\{\,\bigr\}\,\bigl\lfloor
  \,\bigr\rfloor\,\bigl\lceil\,\bigr\rceil\,\bigl\langle\,\bigr\rangle
  \,\big/\,\big\backslash\,\big\vert\,\big\Vert\,\bigm\uparrow\,\bigm\Uparrow
  \,\bigm\downarrow\,\bigm\Downarrow\,\bigm\updownarrow\,\bigm\Updownarrow$\cr
\enddisplay
You can also type ^|\Bigl| and ^|\Bigr| to get larger symbols suitable for
displays:
\begindisplay
$\Bigl(\,\Bigr)\,\Bigl[\,\Bigr]\,\Bigl\{\,\Bigr\}\,\Bigl\lfloor
  \,\Bigr\rfloor\,\Bigl\lceil\,\Bigr\rceil\,\Bigl\langle\,\Bigr\rangle
  \,\Big/\,\Big\backslash\,\Big\vert\,\Big\Vert\,\Bigm\uparrow\,\Bigm\Uparrow
  \,\Bigm\downarrow\,\Bigm\Downarrow\,\Bigm\updownarrow\,\Bigm\Updownarrow$
\enddisplay
These are 50\% taller than their |\big| counterparts. Displayed formulas
most often use delimiters that are even taller (twice the size of\/ |\big|);
such delimiters are constructed by ^|\biggl| and ^|\biggr|, and they
look like this:
\begindisplay
$\biggl(\,\biggr)\,\biggl[\,\biggr]\,\biggl\{\,\biggr\}\,\biggl\lfloor
  \,\biggr\rfloor\,\biggl\lceil\,\biggr\rceil\,\biggl\langle\,\biggr\rangle
  \,\bigg/\,\bigg\backslash\,\bigg\vert\,\bigg\Vert\,\biggm\uparrow
  \,\biggm\Uparrow\,\biggm\downarrow\,\biggm\Downarrow\,\biggm\updownarrow
  \,\biggm\Updownarrow$
\enddisplay
Finally, there are ^|\Biggl| and ^|\Biggr| versions, 2.5 times as tall
as the |\bigl| and |\bigr| delimiters:
\begindisplay
$\Biggl(\,\Biggr)\,\Biggl[\,\Biggr]\,\Biggl\{\,\Biggr\}\,\Biggl\lfloor
  \,\Biggr\rfloor\,\Biggl\lceil\,\Biggr\rceil\,\Biggl\langle\,\Biggr\rangle
  \,\Bigg/\,\Bigg\backslash\,\Bigg\vert\,\Bigg\Vert\,\Biggm\uparrow
  \,\Biggm\Uparrow\,\Biggm\downarrow\,\Biggm\Downarrow\,\Biggm\updownarrow
  \,\Biggm\Updownarrow$
\enddisplay

\medskip
\exercise Guess how to type the formula $\displaystyle
\biggl({\partial^2\over\partial x^2}+{\partial^2\over\partial y^2}
  \biggr)\bigl\vert\varphi(x+iy)\bigr\vert^2=0$, in display style,
using |\bigg| delimiters for the large parentheses. \ (The symbols $\partial$
and $\varphi$ that appear here are called ^|\partial| and ^|\varphi|.)
\answer |$\displaystyle\biggl({\partial^2\over\partial x^2}+|\hfil\break
|{\partial^2\over\partial y^2}\biggr)\bigl|\||\varphi(x+iy)\bigr|\||^2=0$|.

\dangerexercise In practice, |\big| and |\bigg| delimiters are used much
more often than |\Big| and |\Bigg| ones. Why do you think this is true?
\answer Formulas that are more than one line tall are usually two lines tall,
not 1$1\over2$ or 2$1\over2$ lines tall.

\danger A |\bigl| or |\Bigl| or |\biggl| or |\Biggl| delimiter is an
^{opening}, like a left parenthesis;
a |\bigr| or |\Bigr| or |\biggr| or |\Biggr| delimiter is a
^{closing}, like a right parenthesis. Plain \TeX\ also provides
^|\bigm| and ^|\Bigm| and ^|\biggm| and ^|\Biggm| delimiters, for use
in the middle of formulas; such a delimiter plays the r\^ole of a ^{relation},
like an equals sign, so \TeX\ puts a bit of space on either side of it.
\beginlongmathdemo
|$\bigl(x\in A(n)\bigm|\||x\in B(n)\bigr)$|&
  \tenmath\bigl(x\in A(n)\bigm\vert x\in B(n)\bigr)\cr
\noalign{\vskip2pt}
|$\bigcup_n X_n\bigm\|\||\bigcap_n Y_n$|&
  \tenmath\bigcup_n X_n\bigm\Vert\bigcap_n Y_n\cr
\endmathdemo
^^|\bigcup|^^|\bigcap|^^|\verticalline|^^|\in|
You can also say just ^|\big| or ^|\Big| or ^|\bigg| or ^|\Bigg|; this produces
a delimiter that acts as an ordinary variable. It is used primarily with
slashes and backslashes, as in the following example.
\beginlongmathdemo
\noalign{\vskip-2pt}
|$${a+1\over b}\bigg/{c+1\over d}$$|&
  \tenmath\displaystyle{a+1\over b}\bigg/{c+1\over d}\cr
\endmathdemo

\dangerexercise What's the professional way to type
$\tenmath\bigl(x+f(x)\bigr)\big/\bigl(x-f(x)\bigr)$? \ (Look closely.)
\answer |$\bigl(x+f(x)\bigr) \big/ \bigl(x-f(x)\bigr)$|. \ Notice especially
the `|\big/|'; an ordinary ^{slash} would look too small between the
|\big| parentheses.

\TeX\ has a built-in mechanism that figures out how tall a pair of delimiters
needs to be, in order to enclose a given subformula; so you can use this
method, instead of deciding whether a delimiter should be |\big| or
|\bigg| or whatever. All you do is say
\begindisplay
^|\left|\<delim$_1$>\<subformula>^|\right|\<delim$_2$>
\enddisplay
and \TeX\ will typeset the subformula, putting the specified delimiters at
the left and the right. The size of the delimiters will be just big enough
to cover the subformula. For example, in the display
\beginlongdisplaymathdemo
|$$1+\left(1\over1-x^2\right)^3$$|&1+\left(1\over1-x^2\right)^3\cr
\endmathdemo
\TeX\ has chosen |\biggl(| and |\biggr)|, because smaller delimiters
would be too small for this particular fraction. A simple formula like
`|$\left(x\right)$|' yields just `$\left(x\right)$'; thus, |\left| and
|\right| sometimes choose delimiters that are smaller than |\bigl| and |\bigr|.

Whenever you use |\left| and |\right| they must pair up with each other,
just as braces do in groups. You can't have |\left| in one formula
and |\right| in another, nor are you allowed to type things like
`|\left(...{...\right)...}|' or
`|\left(...\begingroup...\right)...\endgroup|'.
This restriction makes sense, because \TeX\ needs to typeset the
subformula that appears between |\left| and |\right| before it can decide
how big to make the delimiters.  But it is worth explicit mention here,
because you do {\sl not\/} have to match ^{parentheses} and ^{brackets}, etc.,
^^{crotchets, see brackets}
when you are not using |\left| and |\right|: \TeX\ will not complain if
you input a formula like `|$[0,1)$|' or even `|$)($|' or just `|$)$|'.\
(And it's a good thing \TeX\ doesn't, for such unbalanced formulas occur
surprisingly often in mathematics papers.) \ Even when you do use |\left|
and |\right|, \TeX\ doesn't look closely at the particular delimiters that
you happen to choose; thus, you can type strange things like `|\left)|'
and/or `|\right(|' if you know what you're doing. Or even if you don't.

The |\over| operation in the example displayed above does not involve the
`|1+|' at the beginning of the formula; this happens because |\left| and
|\right| have the function of ^{grouping}, in addition to their function
of delimiter-making.  Any definitions that you happen to make between
|\left| and |\right| will be local, as if braces had appeared around the
enclosed subformula.

\exercise Use |\left| and |\right| to typeset the following display
(with ^|\phi| for $\phi$):
$$\pi(n)=\sum_{k=2}^n\left\lfloor\phi(k)\over k-1\right\rfloor.$$
\answer |$$\pi(n)=\sum_{k=2}^n\left\lfloor\phi(k)\over k-1\right\rfloor.$$|

At this point you are probably wondering why you should bother learning about
|\bigl| and |\bigr| and their relatives, when |\left| and |\right| are there
to calculate sizes for you automatically. Well, it's true that |\left|
and |\right| are quite handy, but there are at least three situations in which
you will want to use your own wisdom when selecting the proper delimiter size:
\ (1)~Sometimes |\left| and |\right| choose a smaller delimiter than you want.
For example, we used |\bigl| and |\bigr| to produce $\bigl\vert\vert x\vert-
\vert y\vert\bigr\vert$ in one of the previous illustrations; |\left| and
|\right| don't make things any bigger than necessary, so
`|$\left|\||\left|\||x\right|\||-\left|\||y\right|\||\right|\||$|'
yields only
`$\left\vert \left\vert x\right\vert -\left\vert y\right\vert \right\vert$'.
\ (2)~Sometimes |\left| and |\right| choose a larger delimiter than you want.
This happens most frequently when they enclose a large operator in a display;
for example, compare the following two formulas:
\beginlongdisplaymathdemo
\noalign{\vskip 6pt}
|$$\left( \sum_{k=1}^n A_k \right)$$|&\left( \sum_{k=1}^n A_k \right)\cr
\noalign{\vskip 3pt}
|$$\biggl( \sum_{k=1}^n A_k \biggr)$$|&\biggl( \sum_{k=1}^n A_k \biggr)\cr
\endmathdemo
The rules of\/ |\left| and |\right| cause them to enclose the ^|\sum| together
with its ^{limits}, but in special cases like this it looks better to let
the limits hang out a~bit; |\bigg| delimiters are better here.
\ (3)~Sometimes you need to break a huge displayed
formula into two or more separate lines, and you want to make sure that
its opening and closing delimiters have the same size; but you can't use
|\left| on the first line and |\right| on the last, since |\left| and
|\right| must occur in pairs. The solution is to use |\Biggl| (say) on
the first line and |\Biggr| on the last.

\danger Of course, one of the advantages of\/ |\left| and |\right| is that
they can make arbitrarily large delimiters---much bigger than |\biggggg|!
The slashes and angle brackets do have a maximum size, however; if you
ask for really big versions of those symbols you will get the largest
ones available.

\exercise Prove that you have mastered delimiters: Coerce \TeX\ into
producing the formula
$$\pi(n)=\sum_{m=2}^n\left\lfloor\biggl(\sum_{k=1}^{m-1}\bigl\lfloor
  (m/k)\big/\lceil m/k\rceil\bigr\rfloor\biggr)^{-1}\right\rfloor.$$
\answer |$$\pi(n)=\sum_{m=2}^n\left\lfloor\biggl(\sum_{k=1}^{m-1}\bigl|
\hfil\break
|\lfloor(m/k)\big/\lceil m/k\rceil\bigr\rfloor\biggr)^{-1}\right\rfloor.$$|

\danger If you type `|.|'\ after |\left| or |\right|, instead of
specifying one of the basic delimiters, you get a so-called ^{null
delimiter} (which is blank). Why on earth would anybody want that, you may
ask. Well, you sometimes need to produce formulas that contain only one
large delimiter. For example, the display
$$\vert x\vert=\cases{x,&if $x\ge0$\cr
  -x,&if $x<0$\cr}$$
has a `$\{$' but no `$\}$'. It can be produced by a construction of the form
\begindisplay
|$$|\||x|\||=\left\{ ... \right.$$|
\enddisplay
Chapter 18 explains how to fill in the `\hbox{|...|}' to finish this
construction; let's just notice for now that the `|\right.|'\ makes it
possible to have an invisible right delimiter to go with the visible
left brace.

\ddanger A null delimiter isn't completely void; it is an empty box
whose width is a \TeX\ parameter called ^|\nulldelimiterspace|.
We will see later that null delimiters are inserted next to fractions.
Plain \TeX\ sets |\nulldelimiterspace=1.2pt|.

You can type `|<|' or `|>|' as convenient abbreviations for ^|\langle| and
^|\rangle|, when \TeX\ is looking for a delimiter. For example,
`|\bigl<|' is equivalent to `|\bigl\langle|', and `|\right>|' is
equivalent to `|\right\rangle|'. Of course `|<|' and `|>|' ordinarily
produce the ^{less-than} and ^{greater-than} relations `${<}\,{>}$', which
are quite different from ^{angle brackets} `$\langle\,\rangle$'.

\danger Plain \TeX\ also makes available a few more delimiters, which were
not listed in the basic set of~22 because they are sort of special.
The control sequences ^|\arrowvert|, ^|\Arrowvert|, and ^|\bracevert| produce
delimiters made from the repeatable parts of the vertical arrows, double
vertical arrows, and large braces, respectively, without the arrowheads
or the curly parts of the braces. They produce results similar to
^|\vert| or ^|\Vert|, but they are surrounded by more white space and
they have a different weight. You can also use ^|\lgroup| and ^|\rgroup|,
which are constructed from braces without the middle parts; and
^|\lmoustache| and ^|\rmoustache|, ^^{moustaches}
which give you the top and bottom halves of large braces. For example,
here are the |\Big| and |\bigg| versions of\/ |\vert|, |\Vert|, and
these seven special delimiters:
$$\halign{\indent$#\hfil$\cr
\ldots\Big\vert\ldots\Big\Vert
\ldots\Big\arrowvert\ldots\Big\Arrowvert\ldots\Big\bracevert
\ldots\Big\lgroup\ldots\Big\rgroup\ldots\Big\lmoustache\ldots\Big\rmoustache
\ldots\,;\cr
\noalign{\smallskip}
\ldots\bigg\vert\ldots\bigg\Vert
\ldots\bigg\arrowvert\ldots\bigg\Arrowvert\ldots\bigg\bracevert
\ldots\bigg\lgroup\ldots\bigg\rgroup\ldots\bigg\lmoustache\ldots\bigg\rmoustache
\ldots\,.\cr}$$
Notice that |\lgroup| and |\rgroup| are rather like bold parentheses, with
sharper bends at the corners; this makes them attractive for certain large
displays. But you cannot use them exactly like parentheses, because
they are available only in large sizes (|\Big|~or~more).

\ddanger Question: What happens if a ^{subscript} or ^{superscript}
follows a large delimiter? Answer:~That's a good question. After a |\left|
delimiter, it is the first subscript or superscript of the enclosed
subformula, so it is effectively preceded by |{}|. After a |\right|
delimiter, it is a subscript or superscript of the entire |\left...\right|
subformula. And after a |\bigl| or |\bigr| or |\bigm| or |\big| delimiter,
it applies only to that particular delimiter. Thus, `|\bigl(_2|' works
quite differently from `|\left(_2|'.

\danger If you look closely at the examples of math typesetting in this
chapter, you will notice that large parentheses and brackets are
symmetric with respect to an invisible horizontal line that runs a little
bit above the ^{baseline}; when a delimiter gets larger, its height and
depth both grow by the same amount. This horizontal line is called the
{\sl^{axis}\/} of the formula; for example, a formula in the text of the
present paragraph would have an axis at this level: $\hskip 2em\over$. The
bar line in every fraction is centered on the axis, regardless of the size
of the numerator or denominator.

\danger Sometimes it is necessary to create a special box that should be
centered vertically with respect to the axis. \ (For example, the
`$\vert x\vert=\bigl\{\,\ldots$' example above was done with such a box.) \
\TeX\ provides a simple way to do this: You just say
\begindisplay
|\vcenter{|\<vertical mode material>|}|
\enddisplay
and the vertical mode material will be packed into a box just as if
^|\vcenter| had been ^|\vbox|. Then the box will be raised or lowered until
its top edge is as far above the axis as the bottom edge is below.

\ddanger The concept of ``axis'' is meaningful for \TeX\ only in math
formulas, not in ordinary text; therefore \TeX\ allows you to use
|\vcenter| only in math mode. If you really need to center something
vertically in horizontal mode, the solution is to say `|$\vcenter{...}$|'.
 \ (Incidentally, the constructions `|\vcenter| |to|\<dimen>'
and `|\vcenter| |spread|\<dimen>' are legal too, in math mode;
vertical glue is always set by the rules for |\vbox| in
Chapter~12. But |\vcenter| by itself is usually sufficient.)

\danger Any box can be put into a formula by simply saying ^|\hbox| or
|\vbox| or ^|\vtop| or ^|\box| or ^|\copy| in the normal way, even when
you are in math mode.  Furthermore you can use ^|\raise| or ^|\lower|, as
if you were in horizontal mode, and you can insert vertical rules with
^|\vrule|.  Such constructions, like |\vcenter|, produce boxes that can be
used like ordinary symbols in math formulas.

\ddanger Sometimes you need to make up your own symbols, when you run across
something unusual that doesn't occur in the fonts. If the new symbol
occurs only in one place, you can use |\hbox| or |\vcenter| or something
to insert exactly what you want; but if you are defining a macro for
general use, you may want to use different constructions in different
styles. \TeX\ has a special feature called ^|\mathchoice| that comes
to the rescue in such situations: You write
\begindisplay
|\mathchoice{|\<math>|}{|\<math>|}{|\<math>|}{|\<math>|}|
\enddisplay
where each \<math> specifies a subformula. \TeX\ will choose the first
subformula in style $D$ or~$D'$, the second in style $T$ or~$T'$, the
third in style $S$ or~$S'$, the fourth in style $\SS$ or $\SS'$.
\ (\TeX\ actually typesets all four subformulas, before it chooses the
final one, because the actual style is not always known at the time a
|\mathchoice| is encountered; for example, when you type `|\over|' you often
change the style of everything that has occurred earlier in the formula.
Therefore |\mathchoice| is somewhat expensive in terms of time and space,
and you should use it only when you're willing to pay the price.)

\ddangerexercise Guess what output is produced by the following commands:
\begintt
\def\puzzle{{\mathchoice{D}{T}{S}{SS}}}
$$\puzzle{\puzzle\over\puzzle^{\puzzle^\puzzle}}$$
\endtt
\answer A displayed formula equivalent to |$${D}{{T}\over{T}^{{S}^{SS}}}$$|.

\ddangerexercise Devise a `^|\square|' macro that produces a
\def\sqr#1#2{{\vcenter{\vbox{\hrule height.#2pt
      \hbox{\vrule width.#2pt height#1pt \kern#1pt \vrule}
      \hrule height.#2pt}}}}%
`$\,\sqr34\,$' for use in math formulas. The box should be symmetrical
with respect to the axis, and its inside dimensions should be $3\pt$ in
display and text styles, $2.1\pt$ in script styles, and $1.5\pt$ in
scriptscript styles. The rules should be $0.4\pt$ thick in
display and text styles, $0.3\pt$ thick otherwise.
\answer |\def\sqr#1#2{{\vcenter{\vbox{\hrule height.#2pt|\parbreak
        |        \hbox{\vrule width.#2pt height#1pt \kern#1pt|\parbreak
        |           \vrule width.#2pt}|\parbreak
        |        \hrule height.#2pt}}}}|\parbreak
        |\def\square{\mathchoice\sqr34\sqr34\sqr{2.1}3\sqr{1.5}3}|

\ddanger Plain \TeX\ has a macro called ^|\mathpalette| that is useful
for |\mathchoice| constructions; `|\mathpalette\a{xyz}|' expands to
the four-pronged array of choices
`|\mathchoice|\stretch|{\a|\stretch|\displaystyle|\stretch|{xyz}}|\stretch
|...|\stretch|{\a|\stretch|\scriptscriptstyle|\stretch|{xyz}}|\stretch'.
Thus the first argument to |\mathpalette| is a control sequence whose
first argument is a style selection. Appendix~B contains several examples
that show how |\mathpalette| can be applied. \ (See in particular the
definitions of\/ |\phantom|, |\root|, and |\smash|; the ^{congruence sign}
^|\cong| ($\cong$) is also constructed from $=$ and $\sim$ using
|\mathpalette|.)
^^{constructing new math symbols}
^^{math symbols, construction of}

\ddanger At the beginning of this chapter we discussed the commands
|\over|, |\atop|, |\choose|, and |\above|. These are special cases of
\TeX's ``^{generalized fraction}'' feature, which includes also the
three primitives
\begindisplay
|\overwithdelims|\<delim$_1$>\<delim$_2$>\cr
|\atopwithdelims|\<delim$_1$>\<delim$_2$>\cr
|\abovewithdelims|\<delim$_1$>\<delim$_2$>\<dimen>\cr
\enddisplay
The third of these is the most general, as it encompasses all of the other
generalized fractions:  ^|\overwithdelims| uses a ^{fraction} bar whose
thickness is the default for the current size, and ^|\atopwithdelims| uses
an invisible fraction bar whose thickness is zero, while
^|\abovewithdelims| uses a bar whose thickness is specified explicitly.
\TeX\ places the immediately preceding subformula (the ^{numerator}) over
the immediately following subformula (the ^{denominator}), separated by a
bar line of the desired thickness; then it puts \<delim$_1$> at the left
and \<delim$_2$> at the right. For example, `^|\choose|' is equivalent to
`|\atopwithdelims()|'. If you define |\legendre| to be
`|\overwithdelims()|', you can typeset the ^{Legendre symbol}
\def\legendre{\overwithdelims()}%
`$a\legendre b$' by saying `|{a\legendre b}|'. The size of the surrounding
delimiters depends only on the style, not on the size of the fractions;
larger delimiters are used in styles $D$ and~$D'$ (see Appendix~G\null). The
simple commands ^|\over|, ^|\atop|, and ^|\above| are equivalent to the
corresponding `|withdelims|' commands when the delimiters are null; for
example, `|\over|' is an abbreviation for `|\overwithdelims..|'.

\def\euler{\atopwithdelims<>}
\ddangerexercise Define a control sequence |\euler| so that the
^{Eulerian number} $n\euler k$ will be produced when you type `|{n\euler k}|'
in a formula.
\answer|\def\euler{\atopwithdelims<>}|.

\ddanger Appendix G explains exactly how \TeX\ computes the desired size
of delimiters for |\left| and~|\right|. The general idea is that delimiters
are vertically centered with respect to the ^{axis}; hence, if we want
to cover a subformula between |\left| and |\right| that extends $y_1$~units
above the axis and $y_2$~units below, we need to make a delimiter whose
height plus depth is at least $y$~units, where $y=2\max(y_1,y_2)$.
It is usually best not to cover the formula completely, however,
but just to come close; so \TeX\ allows you to specify
two parameters, the ^|\delimiterfactor|~$f$ (an~integer) and the
^|\delimitershortfall|~$\delta$ (a~dimension). The minimum delimiter size
is taken to be at least $y\cdot f/1000$, and at least $y-\delta$. Appendix~B
sets $f=901$ and $\delta=5\pt$. Thus, if $y=30\pt$, the plain \TeX\ format
causes the delimiter to be more than $27\pt$ tall; if $y=100\pt$, the
corresponding delimiter will be at least $95\pt$ tall.

\danger So far we have been discussing the rules for typing math formulas,
but we haven't said much about how \TeX\ actually goes about converting
its input into lists of boxes and glue. Almost all of the control
sequences that have been mentioned in Chapters 16 and~17 are ``high level''
features of the plain \TeX\ format; they are not built into \TeX\ itself.
Appendix~B defines those control sequences in terms of more primitive
commands that \TeX\ actually deals with. For example, `|\choose|' is
an abbreviation for `|\atopwithdelims()|'; Appendix~B not only introduces
|\choose|, it also tells \TeX\ where to find the delimiters |(| and~|)|
in various sizes. The plain \TeX\ format defines all of the special
characters like |\alpha| and~|\mapsto|, all of the special accents like
|\tilde| and~|\widehat|, all of the large operators like |\sum| and~|\int|,
and all of the delimiters like |\lfloor| and~|\vert|. Any of these things
can be redefined, in order to adapt \TeX\ to other mathematical styles
and/or to other fonts.

\danger The remainder of this chapter discusses the low-level commands
that \TeX\ actually obeys behind the scenes.  Every paragraph on the next
few pages is marked with double dangerous bends, so you should skip to
Chapter~18 unless you are a glutton for \TeX nicalities.

\ninepoint
\ddanger All characters that are typeset in math mode belong to one of
sixteen {\sl^{families} of fonts}, numbered internally from 0 to~15. Each
of these families consists of three fonts: one for text size, one for
script size, and one for scriptscript size. The commands ^|\textfont|,
^|\scriptfont|, and ^|\scriptscriptfont| are used to specify the members
of each family.  For example, ^{family~0} in the plain \TeX\ format is
used for roman letters, and Appendix~B contains the instructions
\begintt
\textfont0=\tenrm
\scriptfont0=\sevenrm
\scriptscriptfont0=\fiverm
\endtt
to set up this family: The 10-point roman font (^|\tenrm|) is used for
normal symbols, 7-point roman (^|\sevenrm|) is used for subscripts, and
5-point roman (^|\fiverm|) is used for sub-subscripts. Since there are up to
256~characters per font, and 3~fonts per family, and 16~families, \TeX\ can
access up to 12,288 characters in any one formula (4096 in~each of
the three sizes). Imagine that.

\ddanger A definition like |\textfont|\<family number>|=|\<font identifier>
is local to the group that contains it, so you can easily change family
membership from one set of conventions to another and back again. Furthermore
you can put any font into any family; for example, the command
\begintt
\scriptscriptfont0=\scriptfont0
\endtt
makes sub-subscripts in family~0 the same size as the subscripts currently
are. \TeX\ doesn't check to see if the families are sensibly organized; it
just follows instructions. \ (However, fonts cannot be used in families
2 and~3 unless they contain a certain number of special parameters, as we
shall see later.) \ Incidentally, \TeX\ uses ^|\nullfont|, which contains
no characters, for each family member that has not been defined.

\ddanger During the time that a math formula is being read,
\TeX\ remembers each symbol as being ``character position so-and-so in
family number such-and-such,'' but it does not take note of what fonts
are actually in the families until reaching the end of the formula.
Thus, if you have loaded a font called |\Helvetica| that contains Swiss-style
numerals, and if you say something like
\begintt
$\textfont0=\tenrm 9 \textfont0=\Helvetica 9$
\endtt
you will get two 9's in font |\Helvetica|, assuming that \TeX\ has been
set up to take 9's from family~0. The reason is that |\textfont0|
is~|\Helvetica| at the end of the formula, and that's when it counts. On
the other hand, if you say
\begintt
$\textfont0=\tenrm 9 \hbox{$9\textfont0=\Helvetica$}$
\endtt
the first 9 will be from |\tenrm| and the second from |\Helvetica|, because
the formula in the hbox will be typeset before it is incorporated into
the surrounding formula.

\ddangerexercise If you say `|${\textfont0=\Helvetica 9}$|', what
font will be used for the~9?
\answer The |\textfont0| that was current at the beginning of the formula
will be used, because this redefinition is local to the braces. \
(It would be a different story if `^|\global||\textfont|' had appeared instead;
that would have changed the meaning of\/ |\textfont0| at all levels.)

\ddanger Every ^{math character} is given an identifying code
number between 0 and~4095, obtained by adding 256~times the family number
to the position number. This is easily expressed in ^{hexadecimal
notation}, using one hexadecimal digit for the family and two for the
character; for example, \hex{24A} stands for character~\hex{4A} in
family~2. Each character is also assigned to one of eight classes,
^^{classes of math characters, table} ^^{math codes} ^^{table of ...}
numbered 0 to~7, as follows:
$$\halign{\indent#\hfil&\quad#\hfil&\quad#\hfil&
\hskip4em#\hfil&\quad#\hfil&\quad#\hfil\cr
\it \kern-2pt Class&\it Meaning&\kern-2pt\it Example&
\it \kern-2pt Class&\it Meaning&\kern-2pt\it Example\cr
\noalign{\vskip2pt}
0&Ordinary&|/|&
4&Opening&|(|\cr
1&Large operator&|\sum|&
5&Closing&|)|\cr
2&Binary operation&|+|&
6&Punctuation&|,|\cr
3&Relation&|=|&
7&Variable family&|x|\cr
}$$
^^{large operator}^^{binary operation}^^{relation}^^{opening}^^{closing}
^^{punctuation}^^{variable family}
Classes 0 to 6 tell what ``part of speech'' the character belongs to, in
math-printing language; class~7 is a special case discussed below. The class
number is multiplied by 4096 and added to the character number, and this
is the same as making it the leading digit of a four-digit hexadecimal
number. For example, Appendix~B defines |\sum| to be the math character
\hex{1350}, meaning that it is a large operator (class~1) found in position
\hex{50} of family~3.

\ddangerexercise The ^|\oplus| and ^|\bullet| symbols ($\oplus$ and $\bullet$)
are binary operations that appear in positions 8 and~15 (decimal)
of family~2, when the fonts of plain~\TeX\ are being used. Guess
what their math character codes are. \ (This is too easy.)
\answer \hex{2208} and \hex{220F}.

\ddanger Class 7 is a special case that allows math symbols to change families.
It behaves exactly like class~0, except that the specified family is
replaced by the current value of an integer parameter called ^|\fam|,
provided that |\fam| is a legal family number (i.e., if it lies between
0 and~15). \TeX\ automatically sets |\fam=-1| whenever math mode is entered;
therefore class~7 and class~0 are equivalent unless |\fam| has been
given a new value. Plain \TeX\ changes |\fam| to~0 when the user
types `^|\rm|'; this makes it convenient to get roman letters in formulas,
as we will see in Chapter~18, since letters belong to class~7. \ (The
control sequence |\rm| is an abbreviation for `|\fam=0 \tenrm|'; thus,
|\rm| causes
|\fam| to become zero, and it makes |\tenrm| the ``^{current font}.''
In horizontal mode, the |\fam| value is irrelevant and the current font
governs the typesetting of letters; but in math mode, the current font is
irrelevant and the |\fam| value governs the letters. The current font
affects math mode only if\/ |\|\] is used ^^{control space} or if
dimensions are given in ^|ex| or ^|em| units;
it also has an effect if an |\hbox| appears inside a formula, since
the contents of an hbox are typeset in horizontal mode.)

\ddanger The interpretation of characters in math mode is defined by a
table of~256 ``mathcode'' values; these table entries can be changed
by the ^|\mathcode| command, just as the category codes are changed
by ^|\catcode| (see Chapter~7). Each mathcode specifies class, family, and
character position, as described above. For example, Appendix~B contains
the commands
\begintt
\mathcode`<="313C
\mathcode`*="2203
\endtt
which cause \TeX\ to treat the character `|<|' in math mode as a relation
^^{less than} (class~3) found in position \hex{3C} of family~1, and to treat an
^{asterisk} `|*|' as a binary operation found in position~3 of family~2.
The initial value of\/ |\mathcode`b| is \hex{7162}; thus, |b|~is character
\hex{62} in ^{family~1} (italics), and its family will vary with |\fam|.
\ (|INITEX| starts out with |\mathcode|$\,x=x$ for all characters~$x$
that are neither
letters nor digits. The ten digits have |\mathcode|$\,x=x+\hbox{\hex{7000}}$;
the 52 letters have |\mathcode|$\,x=x+\hbox{\hex{7100}}$.) \
\TeX\ looks at the mathcode only when it is typesetting a character whose
catcode is 11~(letter) or 12~(other), or when it encounters a character that
is given explicitly as ^|\char|\<number>.

\ddanger A |\mathcode| can also have the special value \hex{8000}, which
causes the character to behave as if it has catcode~13 (active). Appendix~B
uses this feature to make |'| ^^{apostrophe} expand to |^{|^|\prime||}| in a
slightly tricky way. The mathcode of |'| does not ^^{active math character}
interfere with the use of |'| in ^{octal} constants.

\ddanger The mathcode table allows you to refer indirectly to any character in
any family, with the touch of a single key. You can also specify a math
character code directly, by typing ^|\mathchar|, which is analogous to
^|\char|. For example, the command `|\mathchar"1ABC|' specifies a
character of class~1, family~10 (\hex A), and position \hex{BC}. A~hundred
or so definitions like
\begintt
\def\sum{\mathchar"1350 }
\endtt
would therefore suffice to define the special symbols of plain \TeX\null. But
there is a better way: \TeX\ has a primitive command ^|\mathchardef|,
which relates to |\mathchar| just as ^|\chardef| does to |\char|.
Appendix~B has a hundred or so definitions like
\begintt
\mathchardef\sum="1350
\endtt
to define the special symbols. A |\mathchar| must be between 0 and 32767
(\hex{7FFF}).

\ddanger A character of class~1, i.e., a ^{large operator} like |\sum|, will
be vertically centered with respect to the axis when it is typeset. Thus,
the large operators can be used with different sizes of type. This vertical
adjustment is not made for symbols of the other classes.

\ddanger \TeX\ associates classes with subformulas as well as with individual
characters. Thus, for example, you can treat a complex construction as if
it were a binary operation or a relation, etc., if you want to. The
commands ^|\mathord|, ^|\mathop|, ^|\mathbin|, ^|\mathrel|, ^|\mathopen|,
^|\mathclose|, and ^|\mathpunct| are used for this purpose; each of them
is followed either by a single character or by a subformula in braces.
For example, |\mathopen\mathchar"1234| is equivalent to |\mathchar"4234|,
because |\mathopen| forces class~4 (opening). In the formula
`|$G\mathbin:H$|', the ^{colon} is~treated as a binary operation.
And Appendix~B constructs large opening symbols by
defining ^|\bigl||#1| to be an abbreviation for
\begintt|let|0|rm
\mathopen{\hbox{$\left#1|<|0strut>\right.$}}
\endtt
There's also an eighth classification, ^|\mathinner|, which is not
normally used for individual symbols; ^|\left|$\,\ldots\,$^|\right|
constructions are treated as ``inner'' subformulas, which means that
they will be surrounded by additional space in certain circumstances.
All other subformulas are generally treated as ordinary symbols,
whether they are formed by |\overline| or |\hbox| or |\vcenter| or
by simply being enclosed in braces. Thus, |\mathord| isn't really
a necessary part of the \TeX\ language; instead of typing
`|$1\mathord,234$|' you can get the same effect from `|$1{,}234$|'.

\ddangerexercise Commands like |\mathchardef\alpha="010B| are used in
Appendix~B to define the lowercase ^{Greek} letters. Suppose that you want
to extend plain \TeX\ by putting ^{boldface math italic} letters
in family~9, analogous to the normal math italic letters in family~1.
\ (Such fonts aren't available in stripped down versions of \TeX, but
let's assume that they exist.) \ Assume that the control sequence
|\bmit| has been defined as an abbreviation for `|\fam=9|'; hence
`|{\bmit b}|' will give a boldface math italic~|b|. What change to the
definition of\/ |\alpha| will make |{\bmit\alpha}| produce a boldface~alpha?
\checkequals\bmiexno\exno
\answer |\mathchardef\alpha="710B|. Incidentally, |{\rm\alpha}| will
then give a spurious result, because character position \hex{0B} of
roman fonts does not contain an alpha; you should warn
your users about what characters they are allowed to type under the
influence of special conventions like ^|\rm|.

\ddanger ^{Delimiters} are specified in a similar but more complicated
way. Each character has not only a~|\catcode| and a~|\mathcode| but also
a~^|\delcode|, which is either negative (for characters that should not
act as delimiters) or less than \hex{1000000}. In other words,
nonnegative delcodes consist of six hexadecimal digits. The first three
digits specify a ``small'' variant of the delimiter, and the last three
specify a ``large'' variant. For example, the command
\begintt
\delcode`x="123456
\endtt
means that if the letter |x| is used as a delimiter, its small variant
is found in position \hex{23} of family~1, and its large variant is found
in position \hex{56} of family~4. If the small or large variant is
given as |000|, however (position~0 of ^{family~0}), that variant is ignored.
\TeX\ looks at the delcode when a character follows ^|\left| or ^|\right|,
or when a character follows one of the ^|withdelims| commands; a
negative delcode leads to an error message, but otherwise \TeX\ finds
a suitable delimiter by first trying the small variant and then
the large. \ (Appendix~G discusses this process in more detail.) \
For example, Appendix~B contains the commands
\begintt
\delcode`(="028300  \delcode`.=0
\endtt
which specify that the small variant of a left parenthesis is found in
position \hex{28} of family~0, and that the large variant is in position~0
of family~3; also, a period has no variants, hence `|\left.|'\ will produce
a ^{null delimiter}. There actually are several different left parenthesis
symbols in family~3; the smallest is in position~0, and the others are
linked together by information that comes with the font. All delcodes
are~$-1$ until they are changed by a |\delcode| command.

\ddangerexercise Appendix~B defines |\delcode`<| so that there is a
shorthand notation for ^{angle brackets}. Why do you think Appendix~B
doesn't go further and define |\delcode`{|?
\answer Assigning |\delcode`{| would not work to allow `|\left{|', because
the brace has category~1 and isn't a legal \<delim>.
Allowing brace delimiters would be a bad idea because it would
mess up other constructions, such as arguments to macros, and
components of alignments. Moreover, a user who
gets away with `|\left{|' is likely to try also `|\bigl{|', which
fails miserably.

\ddanger A delimiter can also be given directly, as `^|\delimiter|\<number>'.
In this case the number can be as high as \hex{7FFFFFF}, i.e., seven
hexadecimal digits; the leading digit specifies a class, from 0 to~7,
as in a |\mathchar|. For example, Appendix~B contains the definition
\begintt
\def\langle{\delimiter"426830A }
\endtt
and this means that ^|\langle| is an opening (class 4) whose small
variant is \hex{268} and whose large variant is \hex{30A}. When |\delimiter|
appears after |\left| or |\right|, the class digit is ignored; but
when |\delimiter| occurs in other contexts, i.e., when \TeX\ isn't
looking for a delimiter, the three rightmost digits are dropped and
the remaining four digits act as a |\mathchar|. For example, the expression
`|$\langle x$|' is treated as if it were `|$\mathchar"4268 x$|'.

\ddangerexercise What goes wrong if you type
`|\bigl\delimiter"426830A|'\thinspace?
\answer Since |\bigl| is defined as a macro with one parameter,
it gets just `|\delimiter|' as the argument. You have to write
`|\bigl{\delimiter"426830A}|' to make this work. On the other hand,
|\left| will balk if the following character is a left brace. Therefore
it's best to have control sequence names for all delimiters.

\ddanger Granted that these numeric conventions for |\mathchar| and
|\delimiter| are not beautiful, they sure do pack a lot of information into
a small space. That's why \TeX\ uses them for low-level definitions inside
formats. Two other low-level primitives also deserve to be mentioned:
^|\radical| and ^|\mathaccent|. Plain \TeX\ makes ^{square root signs}
and math accents available by giving the commands
\begintt
\def\sqrt{\radical"270370 }
\def\widehat{\mathaccent"362 }
\endtt
and several more like them. The idea is that |\radical| is followed by
a delimiter code and |\mathaccent| is followed by a math character code,
so that \TeX\ knows the family and character positions for the symbols
used in radical and accent constructions. Appendix~G gives precise
information about the positioning of these characters. By changing the
definitions, \TeX\ could easily be extended so that it would typeset a
variety of different radical signs and a variety of different accent
signs, if such symbols were available in the fonts.
^^{surd signs, see radical}

\ddanger Plain \TeX\ uses ^{family~1} for math italic letters, ^{family~2} for
ordinary math symbols, and ^{family~3} for large symbols. \TeX\ insists that
^^{math fonts}
the fonts in families 2 and~3 have special ^|\fontdimen| parameters,
which govern mathematical spacing according to the rules in Appendix~G\null;
the ^|cmsy| and ^|cmex| ^{symbol fonts} have these parameters, so
their assignment to families 2 and~3 is almost mandatory. \ (There is, however,
a way to modify the parameters of any font, using the ^|\fontdimen| command.) \
^|INITEX| initializes the mathcodes of all ^{letters} |A| to~|Z| and |a| to~|z|
so that they are symbols of class~7 and family~1; that's why
it is natural to use family~1 for math italics.  Similarly, the digits |0|
to~|9| are class~7 and family~0.  None of the other families
is treated in any special way by \TeX. Thus, for example, plain \TeX\ puts
^{text italic} in family~4, slanted roman in family~5, bold roman in family~6,
and typewriter type in family~7, but any of these numbers could be
switched around. There is a macro ^|\newfam|, analogous to |\newbox|,
that will assign symbolic names to families that aren't already used.

\ddanger When \TeX\ is in horizontal mode, it is making a horizontal list;
in vertical mode, it is making a vertical list. Therefore it should come
as no great surprise that \TeX\ is making a ^{math list} when it is in
^{math mode}. The contents of horizontal lists were explained in Chapter~14,
and the contents of vertical lists were explained in Chapter~15; it's time
now to describe what math lists are made of. Each item in a math list
is one of the following types of things:\enddanger

\smallskip
\item\bull an ^{atom} (to be explained momentarily);

\item\bull horizontal material (a rule or discretionary or penalty or
``whatsit'');

\item\bull vertical material (from |\mark| or |\insert| or |\vadjust|);

\item\bull a glob of ^{glue} (from |\hskip| or |\mskip| or |\nonscript|);

\item\bull a ^{kern} (from |\kern| or |\mkern|);

\item\bull a ^{style change} (from |\displaystyle|, |\textstyle|, etc.);

\item\bull a ^{generalized fraction} (from |\above|, |\over|, etc.);

\item\bull a ^{boundary} (from |\left| or |\right|);

\item\bull a four-way ^{choice} (from ^|\mathchoice|).

\ddanger The most important items are called {\sl atoms}, and they have
three parts: a {\sl^{nucleus}}, a {\sl^{superscript}}, and a {\sl^{subscript}}.
For example, if you type
\begintt
(x_i+y)^{\overline{n+1}}
\endtt
in math mode, you get a math list consisting of five atoms:
$($, $x_i$, $+$, $y$, and~$)^{\overline{n+1}}$. The nuclei of these atoms
are $($, $x$, $+$, $y$, and~$)$; the subscripts are empty except for the
second atom, which has subscript~$i$; the superscripts are empty except for the
last atom, whose superscript is~$\overline{n+1}$. This superscript is
itself a math list consisting of one atom, whose nucleus is~$n+1$; and that
nucleus is a math list consisting of three atoms.

\ddanger There are thirteen kinds of atoms, each of which might act
differently in a formula; for example, `$($' is an Open atom because
^^{atomic types, table}
it comes from an opening. Here is a complete list of the different kinds:
$$\halign{\indent#\hfil&\enskip#\hfil\cr
Ord&is an ordinary atom like `$x$'\thinspace;\cr
Op&is a large operator atom like `$\sum$'\thinspace;\cr
Bin&is a binary operation atom like `$+$'\thinspace;\cr
Rel&is a relation atom like `$=$'\thinspace;\cr
Open&is an opening atom like `$($'\thinspace;\cr
Close&is a closing atom like `$)$'\thinspace;\cr
Punct&is a punctuation atom like `$,$'\thinspace;\cr
Inner&is an inner atom produced by `|\left|$\,\ldots\,$|\right|';\cr
Over&is an overline atom like `$\overline x$'\thinspace;\cr
Under&is an underline atom like `$\underline x$'\thinspace;\cr
Acc&is an accented atom like `$\hat x$'\thinspace;\cr
Rad&is a radical atom like `$\sqrt2$'\thinspace;\cr
Vcent&is a vbox to be centered, produced by |\vcenter|.\cr
}$$

\ddanger An atom's nucleus, superscript, and subscript are called its
{\sl ^{fields}}, and there are four possibilities for each of these fields.
A field can be\enddanger

\smallskip
\item\bull empty;

\item\bull a math symbol (specified by family and position number);

\item\bull a box; or

\item\bull a math list.

\smallskip\noindent
For example, the Close atom $)^{\overline{n+1}}$ considered above has an
empty subscript field; its nucleus is the symbol `$)$', which is
character~\hex{28} of family~0 if the conventions of plain \TeX\ are
in force; and its superscript field is the math list $\overline{n+1}$.
The latter math list consists of an Over atom whose nucleus
is the math list $n+1$; and that math list, in turn, consists of
three atoms of types Ord, Bin, Ord.

\ddanger You can see \TeX's view of a math list by typing ^|\showlists|
in math mode. ^^{internal list format}
For example, after `|$(x_i+y)^{\overline{n+1}}\showlists|' your log
file gets the following curious data:
\begindisplay
|\mathopen|\cr
|.\fam0 (|\cr
|\mathord|\cr
|.\fam1 x|\cr
|_\fam1 i|\cr
\noalign{\penalty-500}
|\mathbin|\cr
|.\fam0 +|\cr
\noalign{\penalty-500}
|\mathord|\cr
|.\fam1 y|\cr
\noalign{\penalty-500}
|\mathclose|\cr
|.\fam0 )|\cr
|^\overline|\cr
|^.\mathord|\cr
|^..\fam1 n|\cr
|^.\mathbin|\cr
|^..\fam0 +|\cr
|^.\mathord|\cr
|^..\fam0 1|\cr
\enddisplay
In our previous experiences with |\showlists| we observed that there can
be boxes within boxes, and that each line in the log file is
prefixed by dots to indicate its position in the hierarchy. Math lists
have a slightly more complex structure; therefore a dot is used to denote
the nucleus of an atom, a~`|^|' is used for the superscript field, and
a~`|_|' is used for the subscript field. Empty fields are not shown. Thus,
for example, the Ord atom~$x_i$ is represented here by three lines
`|\mathord|', `|.\fam1 x|', and `|_\fam1 i|'.

\ddanger Certain kinds of atoms carry additional information besides their
nucleus, subscript, and superscript fields: An Op atom will be marked
`^|\limits|' or `^|\nolimits|' if the normal ^|\displaylimits|
convention has been overridden; a Rad atom contains
a delimiter field to specify what radical sign is to be used; and an Acc atom
contains the family and character codes of the accent symbol.

\ddanger When you say ^|\hbox||{...}| in math mode, an Ord atom is placed
on the current math list, with the hbox as its nucleus.  Similarly,
^|\vcenter||{...}| produces a Vcent atom whose nucleus is a box. But in
most cases the nucleus of an atom will be either a symbol or a math list.
You can experiment with |\showlists| to discover how other things like
fractions and mathchoices are represented internally.

\ddanger Chapter~26 contains complete details of how math lists are
constructed.  As soon as math mode ends (i.e., when the closing `|$|'
occurs), \TeX\ dismantles the current math list and converts it into a
horizontal list.  The rules for this conversion are spelled out in
Appendix~G\null. You can see ``before and after'' representations of such math
typesetting by ending a formula with `|\showlists$\showlists|'; the first
|\showlists| will display the math list, and the second will show the
(possibly complex) horizontal list that is manufactured from it.

\endchapter

The learning time is short. A few minutes gives the general flavor, and
typing a page or two of a paper generally uncovers most of the misconceptions.
\author ^{KERNIGHAN} and ^{CHERRY}, {\sl A System for %
  Typesetting Mathematics\/} (1975)
  % in Communications of the ACM p152

\bigskip

Within a few hours (a few days at most)
a typist with no math or typesetting experience
can be taught to input even the most complex equations.
\author PETER J. ^{BOEHM}, {\sl Software and Hardware Considerations %
  for a\break Technical Typesetting System\/} (1976)
%  in IEEE Transactions on Professional Communication PC-19, pp15--19

\eject
